% ============================================================
% Appendix: Prompts Used in CELPO
% ============================================================
% Add this to your main LaTeX document with:
%   % ============================================================
% Appendix: Prompts Used in CELPO
% ============================================================
% Add this to your main LaTeX document with:
%   % ============================================================
% Appendix: Prompts Used in CELPO
% ============================================================
% Add this to your main LaTeX document with:
%   % ============================================================
% Appendix: Prompts Used in CELPO
% ============================================================
% Add this to your main LaTeX document with:
%   \input{appendix_prompts}
% or copy the content into your appendix section.
% Requires packages: listings, xcolor, amsmath, tcolorbox (optional)
% ============================================================

\appendix
\section{Prompts}\label{app:prompts}

This appendix provides the complete prompt templates used across the different stages of our CELPO framework: inference (examination), teacher correction (hint generation), training, and answer verification.

% ----------------------------------------------------------
\subsection{System Prompt}\label{app:system-prompt}

All inference and training stages share the same system prompt to activate the model's mathematical reasoning capability:

\begin{tcolorbox}[colback=gray!5, colframe=gray!50, title=System Prompt]
\texttt{Please reason step by step and put your final answer within \textbackslash boxed\{\}.}
\end{tcolorbox}

% ----------------------------------------------------------
\subsection{Student Examination Prompt}\label{app:exam-prompt}

During the examination phase, the student model receives a math problem via the standard chat template. The input is structured as:

\begin{tcolorbox}[colback=blue!3, colframe=blue!40, title=Student Exam Input Format]
\begin{verbatim}
<|im_start|>system
Please reason step by step and put your final 
answer within \boxed{}.<|im_end|>
<|im_start|>user
{question}<|im_end|>
<|im_start|>assistant
\end{verbatim}
\end{tcolorbox}

\noindent The model then generates its reasoning and final answer in a single pass.

% ----------------------------------------------------------
\subsection{Hint-Enhanced Examination Prompt}\label{app:hint-exam-prompt}

When evaluating the student model with teacher-provided hints (prefix-forcing), the hint is prepended to the assistant's response as a prefix:

\begin{tcolorbox}[colback=blue!3, colframe=blue!40, title=Hint-Enhanced Exam Input Format]
\begin{verbatim}
<|im_start|>system
Please reason step by step and put your final 
answer within \boxed{}.<|im_end|>
<|im_start|>user
{question}<|im_end|>
<|im_start|>assistant
{hint}
\end{verbatim}
\end{tcolorbox}

\noindent The model is forced to continue generation after the hint prefix, producing the answer portion conditioned on the provided hint.

% ----------------------------------------------------------
\subsection{Teacher Correction Prompt}\label{app:teacher-correct-prompt}

The teacher model receives the following prompt to generate targeted knowledge hints based on the student's incorrect solution:

\begin{tcolorbox}[colback=green!3, colframe=green!40, title=Teacher Correction Prompt, breakable]
\begin{verbatim}
**Role:** Heuristic Logic Mentor & Knowledge Bridge

**Task:**
Your task is to analyze the **Student's Solution** in 
comparison to the **Reference Answer**. Instead of acting 
as a simple grader, you must identify the specific 
**missing logical link** or **knowledge gap** that prevents 
the student from reaching the correct conclusion.

**Goal:**
Provide the student with the necessary "Knowledge Hints" 
so that, based on their existing correct reasoning, they 
can bridge the gap and solve the problem themselves.

**Critical Constraints for Hints:**
1. **Universal Knowledge Only:** Hints must be provided as 
   **general problem-solving methods, formulas, theorems, 
   or definitions** (e.g., "Recall the formula for the area 
   of a circle: S = pi*r^2" or "Use the derivatives of 
   trigonometric functions: (sin x)' = cos x").
2. **No Specific Calculations:** Do not calculate the result 
   for the specific numbers in the problem. Provide the 
   *tool*, not the *answer*.
3. **Targeted:** The hint must directly address the specific 
   step where the student's logic broke down or stopped.
4. **Token Limit:** Each individual hint must be strictly 
   limited to **50 tokens maximum**. Keep descriptions 
   concise and focus purely on the core principle.
5. **Unidentifiable Gap:** If it is impossible to determine 
   if the student lacks specific universal knowledge (e.g., 
   the error is purely a calculation/arithmetic mistake, a 
   typo, the answer is empty/irrelevant, or the reasoning 
   is too unclear to pinpoint a missing theorem), 
   return `null`.

# Problem:
{problem}

# Student's Answer:
{student_answer}

# Reference Answer:
{ref_solution}

**Output Format:**

**Condition A:** If Constraint #5 applies (unable to identify 
a specific general knowledge gap), output strictly:
`null`

**Condition B:** If a knowledge gap is identified, respond in 
the following XML format:
<hints>
[Provide the general formula, theorem, or principle. 
 Use LaTeX for math expressions.]
...
</hints>

**Example:**
* **Student's Error:** Calculated probability as P(A)+P(B) 
  but events were not mutually exclusive.
* **Your Output:**
<hints>
For any two events A and B, the probability of their union 
is P(A ∪ B) = P(A) + P(B) - P(A ∩ B).
</hints>
\end{verbatim}
\end{tcolorbox}

% ----------------------------------------------------------
\subsection{Training Target Format (IRDCL)}\label{app:training-target}

During CELPO training, the model learns to generate both hints and answers. The training target for Mode~B (hint-augmented generation) concatenates the hint and answer:

\begin{tcolorbox}[colback=orange!3, colframe=orange!40, title=Training Target Format (Mode B)]
\begin{verbatim}
{hints}{answer}
\end{verbatim}
\end{tcolorbox}

\noindent where \texttt{\{hints\}} is the teacher-generated knowledge hint and \texttt{\{answer\}} is the correct reasoning chain ending with \verb|\boxed{...}|. The loss is computed separately on the hint tokens (with Forward KL regularization) and the answer tokens (with Reverse KL regularization).

For anchor samples (Mode~A, pure SFT), only the answer portion is used as the training target:

\begin{tcolorbox}[colback=orange!3, colframe=orange!40, title=Training Target Format (Mode A / Anchor)]
\begin{verbatim}
<|im_start|>system
Please reason step by step and put your final 
answer within \boxed{}.<|im_end|>
<|im_start|>user
{question}<|im_end|>
<|im_start|>assistant
{answer}<|im_end|>
\end{verbatim}
\end{tcolorbox}

% ----------------------------------------------------------
\subsection{Answer Verification Prompt}\label{app:verify-prompt}

To verify whether the student's predicted answer matches the gold answer, we use the following grading prompt:

\begin{tcolorbox}[colback=red!3, colframe=red!40, title=Answer Verification (OREAL Correctness Judge), breakable]
\begin{verbatim}
You are a helpful assistant who evaluates the correctness 
and quality of models' outputs.
Please as a grading expert, judge whether the final answers 
given by the candidates below are consistent with the 
standard answers, that is, whether the candidates answered 
correctly.

Here are some evaluation criteria:
1. Please refer to the given standard answer. You don't need 
   to re-generate the answer to the question because the 
   standard answer has been given. You only need to judge 
   whether the candidate's answer is consistent with the 
   standard answer according to the form of the question. 
   Don't try to answer the original question. You can assume 
   that the standard answer is definitely correct.
2. Because the candidate's answer may be different from the 
   standard answer in the form of expression, before making 
   a judgment, please understand the question and the 
   standard answer first, and then judge whether the 
   candidate's answer is correct, but be careful not to try 
   to answer the original question.
3. Some answers may contain multiple items, such as 
   multiple-choice questions, multiple-select questions, 
   fill-in-the-blank questions, etc. As long as the answer 
   is the same as the standard answer, it is enough. For 
   multiple-select questions and multiple-blank 
   fill-in-the-blank questions, the candidate needs to 
   answer all the corresponding options or blanks correctly 
   to be considered correct.
4. Some answers may be expressed in different ways, such as 
   some answers may be a mathematical expression, some 
   answers may be a textual description, as long as the 
   meaning expressed is the same. And some formulas are 
   expressed in different ways, but they are equivalent 
   and correct.
5. If the prediction is given with \boxed{}, please ignore 
   the \boxed{} and only judge whether the candidate's 
   answer is consistent with the standard answer.

Please judge whether the following answers are consistent 
with the standard answer based on the above criteria. 
Grade the predicted answer of this new question as one of:
A: CORRECT
B: INCORRECT
Just return the letters "A" or "B", with no text around it.

Here is your task. Simply reply with either CORRECT, 
INCORRECT. Don't apologize or correct yourself if there 
was a mistake; we are just trying to grade the answer.

<Original Question Begin>:
{question}
<Original Question End>

<Gold Target Begin>:
{gold_answer}
<Gold Target End>

<Predicted Answer Begin>:
{answer}
<Predicted End>

Judging the correctness of candidates' answers:
\end{verbatim}
\end{tcolorbox}

% ----------------------------------------------------------
\subsection{Summary of Prompt Usage}\label{app:prompt-summary}

Table~\ref{tab:prompt-summary} summarizes the prompts and their usage across the CELPO pipeline stages.

\begin{table}[htbp]
\centering
\caption{Summary of prompts used in the CELPO framework.}
\label{tab:prompt-summary}
\small
\begin{tabular}{lll}
\toprule
\textbf{Prompt} & \textbf{Stage} & \textbf{Purpose} \\
\midrule
System Prompt & All stages & Activate step-by-step reasoning \\
Student Exam & Examination & Generate student solutions \\
Hint-Enhanced Exam & Re-examination & Prefix-force with hints \\
Teacher Correction & Hint Generation & Identify knowledge gaps \\
Training Target (Mode B) & IRDCL Training & Learn hint + answer generation \\
Training Target (Mode A) & Anchor Training & Maintain base capability \\
Answer Verification & Evaluation & Judge answer correctness \\
\bottomrule
\end{tabular}
\end{table}

% or copy the content into your appendix section.
% Requires packages: listings, xcolor, amsmath, tcolorbox (optional)
% ============================================================

\appendix
\section{Prompts}\label{app:prompts}

This appendix provides the complete prompt templates used across the different stages of our CELPO framework: inference (examination), teacher correction (hint generation), training, and answer verification.

% ----------------------------------------------------------
\subsection{System Prompt}\label{app:system-prompt}

All inference and training stages share the same system prompt to activate the model's mathematical reasoning capability:

\begin{tcolorbox}[colback=gray!5, colframe=gray!50, title=System Prompt]
\texttt{Please reason step by step and put your final answer within \textbackslash boxed\{\}.}
\end{tcolorbox}

% ----------------------------------------------------------
\subsection{Student Examination Prompt}\label{app:exam-prompt}

During the examination phase, the student model receives a math problem via the standard chat template. The input is structured as:

\begin{tcolorbox}[colback=blue!3, colframe=blue!40, title=Student Exam Input Format]
\begin{verbatim}
<|im_start|>system
Please reason step by step and put your final 
answer within \boxed{}.<|im_end|>
<|im_start|>user
{question}<|im_end|>
<|im_start|>assistant
\end{verbatim}
\end{tcolorbox}

\noindent The model then generates its reasoning and final answer in a single pass.

% ----------------------------------------------------------
\subsection{Hint-Enhanced Examination Prompt}\label{app:hint-exam-prompt}

When evaluating the student model with teacher-provided hints (prefix-forcing), the hint is prepended to the assistant's response as a prefix:

\begin{tcolorbox}[colback=blue!3, colframe=blue!40, title=Hint-Enhanced Exam Input Format]
\begin{verbatim}
<|im_start|>system
Please reason step by step and put your final 
answer within \boxed{}.<|im_end|>
<|im_start|>user
{question}<|im_end|>
<|im_start|>assistant
{hint}
\end{verbatim}
\end{tcolorbox}

\noindent The model is forced to continue generation after the hint prefix, producing the answer portion conditioned on the provided hint.

% ----------------------------------------------------------
\subsection{Teacher Correction Prompt}\label{app:teacher-correct-prompt}

The teacher model receives the following prompt to generate targeted knowledge hints based on the student's incorrect solution:

\begin{tcolorbox}[colback=green!3, colframe=green!40, title=Teacher Correction Prompt, breakable]
\begin{verbatim}
**Role:** Heuristic Logic Mentor & Knowledge Bridge

**Task:**
Your task is to analyze the **Student's Solution** in 
comparison to the **Reference Answer**. Instead of acting 
as a simple grader, you must identify the specific 
**missing logical link** or **knowledge gap** that prevents 
the student from reaching the correct conclusion.

**Goal:**
Provide the student with the necessary "Knowledge Hints" 
so that, based on their existing correct reasoning, they 
can bridge the gap and solve the problem themselves.

**Critical Constraints for Hints:**
1. **Universal Knowledge Only:** Hints must be provided as 
   **general problem-solving methods, formulas, theorems, 
   or definitions** (e.g., "Recall the formula for the area 
   of a circle: S = pi*r^2" or "Use the derivatives of 
   trigonometric functions: (sin x)' = cos x").
2. **No Specific Calculations:** Do not calculate the result 
   for the specific numbers in the problem. Provide the 
   *tool*, not the *answer*.
3. **Targeted:** The hint must directly address the specific 
   step where the student's logic broke down or stopped.
4. **Token Limit:** Each individual hint must be strictly 
   limited to **50 tokens maximum**. Keep descriptions 
   concise and focus purely on the core principle.
5. **Unidentifiable Gap:** If it is impossible to determine 
   if the student lacks specific universal knowledge (e.g., 
   the error is purely a calculation/arithmetic mistake, a 
   typo, the answer is empty/irrelevant, or the reasoning 
   is too unclear to pinpoint a missing theorem), 
   return `null`.

# Problem:
{problem}

# Student's Answer:
{student_answer}

# Reference Answer:
{ref_solution}

**Output Format:**

**Condition A:** If Constraint #5 applies (unable to identify 
a specific general knowledge gap), output strictly:
`null`

**Condition B:** If a knowledge gap is identified, respond in 
the following XML format:
<hints>
[Provide the general formula, theorem, or principle. 
 Use LaTeX for math expressions.]
...
</hints>

**Example:**
* **Student's Error:** Calculated probability as P(A)+P(B) 
  but events were not mutually exclusive.
* **Your Output:**
<hints>
For any two events A and B, the probability of their union 
is P(A ∪ B) = P(A) + P(B) - P(A ∩ B).
</hints>
\end{verbatim}
\end{tcolorbox}

% ----------------------------------------------------------
\subsection{Training Target Format (IRDCL)}\label{app:training-target}

During CELPO training, the model learns to generate both hints and answers. The training target for Mode~B (hint-augmented generation) concatenates the hint and answer:

\begin{tcolorbox}[colback=orange!3, colframe=orange!40, title=Training Target Format (Mode B)]
\begin{verbatim}
{hints}{answer}
\end{verbatim}
\end{tcolorbox}

\noindent where \texttt{\{hints\}} is the teacher-generated knowledge hint and \texttt{\{answer\}} is the correct reasoning chain ending with \verb|\boxed{...}|. The loss is computed separately on the hint tokens (with Forward KL regularization) and the answer tokens (with Reverse KL regularization).

For anchor samples (Mode~A, pure SFT), only the answer portion is used as the training target:

\begin{tcolorbox}[colback=orange!3, colframe=orange!40, title=Training Target Format (Mode A / Anchor)]
\begin{verbatim}
<|im_start|>system
Please reason step by step and put your final 
answer within \boxed{}.<|im_end|>
<|im_start|>user
{question}<|im_end|>
<|im_start|>assistant
{answer}<|im_end|>
\end{verbatim}
\end{tcolorbox}

% ----------------------------------------------------------
\subsection{Answer Verification Prompt}\label{app:verify-prompt}

To verify whether the student's predicted answer matches the gold answer, we use the following grading prompt:

\begin{tcolorbox}[colback=red!3, colframe=red!40, title=Answer Verification (OREAL Correctness Judge), breakable]
\begin{verbatim}
You are a helpful assistant who evaluates the correctness 
and quality of models' outputs.
Please as a grading expert, judge whether the final answers 
given by the candidates below are consistent with the 
standard answers, that is, whether the candidates answered 
correctly.

Here are some evaluation criteria:
1. Please refer to the given standard answer. You don't need 
   to re-generate the answer to the question because the 
   standard answer has been given. You only need to judge 
   whether the candidate's answer is consistent with the 
   standard answer according to the form of the question. 
   Don't try to answer the original question. You can assume 
   that the standard answer is definitely correct.
2. Because the candidate's answer may be different from the 
   standard answer in the form of expression, before making 
   a judgment, please understand the question and the 
   standard answer first, and then judge whether the 
   candidate's answer is correct, but be careful not to try 
   to answer the original question.
3. Some answers may contain multiple items, such as 
   multiple-choice questions, multiple-select questions, 
   fill-in-the-blank questions, etc. As long as the answer 
   is the same as the standard answer, it is enough. For 
   multiple-select questions and multiple-blank 
   fill-in-the-blank questions, the candidate needs to 
   answer all the corresponding options or blanks correctly 
   to be considered correct.
4. Some answers may be expressed in different ways, such as 
   some answers may be a mathematical expression, some 
   answers may be a textual description, as long as the 
   meaning expressed is the same. And some formulas are 
   expressed in different ways, but they are equivalent 
   and correct.
5. If the prediction is given with \boxed{}, please ignore 
   the \boxed{} and only judge whether the candidate's 
   answer is consistent with the standard answer.

Please judge whether the following answers are consistent 
with the standard answer based on the above criteria. 
Grade the predicted answer of this new question as one of:
A: CORRECT
B: INCORRECT
Just return the letters "A" or "B", with no text around it.

Here is your task. Simply reply with either CORRECT, 
INCORRECT. Don't apologize or correct yourself if there 
was a mistake; we are just trying to grade the answer.

<Original Question Begin>:
{question}
<Original Question End>

<Gold Target Begin>:
{gold_answer}
<Gold Target End>

<Predicted Answer Begin>:
{answer}
<Predicted End>

Judging the correctness of candidates' answers:
\end{verbatim}
\end{tcolorbox}

% ----------------------------------------------------------
\subsection{Summary of Prompt Usage}\label{app:prompt-summary}

Table~\ref{tab:prompt-summary} summarizes the prompts and their usage across the CELPO pipeline stages.

\begin{table}[htbp]
\centering
\caption{Summary of prompts used in the CELPO framework.}
\label{tab:prompt-summary}
\small
\begin{tabular}{lll}
\toprule
\textbf{Prompt} & \textbf{Stage} & \textbf{Purpose} \\
\midrule
System Prompt & All stages & Activate step-by-step reasoning \\
Student Exam & Examination & Generate student solutions \\
Hint-Enhanced Exam & Re-examination & Prefix-force with hints \\
Teacher Correction & Hint Generation & Identify knowledge gaps \\
Training Target (Mode B) & IRDCL Training & Learn hint + answer generation \\
Training Target (Mode A) & Anchor Training & Maintain base capability \\
Answer Verification & Evaluation & Judge answer correctness \\
\bottomrule
\end{tabular}
\end{table}

% or copy the content into your appendix section.
% Requires packages: listings, xcolor, amsmath, tcolorbox (optional)
% ============================================================

\appendix
\section{Prompts}\label{app:prompts}

This appendix provides the complete prompt templates used across the different stages of our CELPO framework: inference (examination), teacher correction (hint generation), training, and answer verification.

% ----------------------------------------------------------
\subsection{System Prompt}\label{app:system-prompt}

All inference and training stages share the same system prompt to activate the model's mathematical reasoning capability:

\begin{tcolorbox}[colback=gray!5, colframe=gray!50, title=System Prompt]
\texttt{Please reason step by step and put your final answer within \textbackslash boxed\{\}.}
\end{tcolorbox}

% ----------------------------------------------------------
\subsection{Student Examination Prompt}\label{app:exam-prompt}

During the examination phase, the student model receives a math problem via the standard chat template. The input is structured as:

\begin{tcolorbox}[colback=blue!3, colframe=blue!40, title=Student Exam Input Format]
\begin{verbatim}
<|im_start|>system
Please reason step by step and put your final 
answer within \boxed{}.<|im_end|>
<|im_start|>user
{question}<|im_end|>
<|im_start|>assistant
\end{verbatim}
\end{tcolorbox}

\noindent The model then generates its reasoning and final answer in a single pass.

% ----------------------------------------------------------
\subsection{Hint-Enhanced Examination Prompt}\label{app:hint-exam-prompt}

When evaluating the student model with teacher-provided hints (prefix-forcing), the hint is prepended to the assistant's response as a prefix:

\begin{tcolorbox}[colback=blue!3, colframe=blue!40, title=Hint-Enhanced Exam Input Format]
\begin{verbatim}
<|im_start|>system
Please reason step by step and put your final 
answer within \boxed{}.<|im_end|>
<|im_start|>user
{question}<|im_end|>
<|im_start|>assistant
{hint}
\end{verbatim}
\end{tcolorbox}

\noindent The model is forced to continue generation after the hint prefix, producing the answer portion conditioned on the provided hint.

% ----------------------------------------------------------
\subsection{Teacher Correction Prompt}\label{app:teacher-correct-prompt}

The teacher model receives the following prompt to generate targeted knowledge hints based on the student's incorrect solution:

\begin{tcolorbox}[colback=green!3, colframe=green!40, title=Teacher Correction Prompt, breakable]
\begin{verbatim}
**Role:** Heuristic Logic Mentor & Knowledge Bridge

**Task:**
Your task is to analyze the **Student's Solution** in 
comparison to the **Reference Answer**. Instead of acting 
as a simple grader, you must identify the specific 
**missing logical link** or **knowledge gap** that prevents 
the student from reaching the correct conclusion.

**Goal:**
Provide the student with the necessary "Knowledge Hints" 
so that, based on their existing correct reasoning, they 
can bridge the gap and solve the problem themselves.

**Critical Constraints for Hints:**
1. **Universal Knowledge Only:** Hints must be provided as 
   **general problem-solving methods, formulas, theorems, 
   or definitions** (e.g., "Recall the formula for the area 
   of a circle: S = pi*r^2" or "Use the derivatives of 
   trigonometric functions: (sin x)' = cos x").
2. **No Specific Calculations:** Do not calculate the result 
   for the specific numbers in the problem. Provide the 
   *tool*, not the *answer*.
3. **Targeted:** The hint must directly address the specific 
   step where the student's logic broke down or stopped.
4. **Token Limit:** Each individual hint must be strictly 
   limited to **50 tokens maximum**. Keep descriptions 
   concise and focus purely on the core principle.
5. **Unidentifiable Gap:** If it is impossible to determine 
   if the student lacks specific universal knowledge (e.g., 
   the error is purely a calculation/arithmetic mistake, a 
   typo, the answer is empty/irrelevant, or the reasoning 
   is too unclear to pinpoint a missing theorem), 
   return `null`.

# Problem:
{problem}

# Student's Answer:
{student_answer}

# Reference Answer:
{ref_solution}

**Output Format:**

**Condition A:** If Constraint #5 applies (unable to identify 
a specific general knowledge gap), output strictly:
`null`

**Condition B:** If a knowledge gap is identified, respond in 
the following XML format:
<hints>
[Provide the general formula, theorem, or principle. 
 Use LaTeX for math expressions.]
...
</hints>

**Example:**
* **Student's Error:** Calculated probability as P(A)+P(B) 
  but events were not mutually exclusive.
* **Your Output:**
<hints>
For any two events A and B, the probability of their union 
is P(A ∪ B) = P(A) + P(B) - P(A ∩ B).
</hints>
\end{verbatim}
\end{tcolorbox}

% ----------------------------------------------------------
\subsection{Training Target Format (IRDCL)}\label{app:training-target}

During CELPO training, the model learns to generate both hints and answers. The training target for Mode~B (hint-augmented generation) concatenates the hint and answer:

\begin{tcolorbox}[colback=orange!3, colframe=orange!40, title=Training Target Format (Mode B)]
\begin{verbatim}
{hints}{answer}
\end{verbatim}
\end{tcolorbox}

\noindent where \texttt{\{hints\}} is the teacher-generated knowledge hint and \texttt{\{answer\}} is the correct reasoning chain ending with \verb|\boxed{...}|. The loss is computed separately on the hint tokens (with Forward KL regularization) and the answer tokens (with Reverse KL regularization).

For anchor samples (Mode~A, pure SFT), only the answer portion is used as the training target:

\begin{tcolorbox}[colback=orange!3, colframe=orange!40, title=Training Target Format (Mode A / Anchor)]
\begin{verbatim}
<|im_start|>system
Please reason step by step and put your final 
answer within \boxed{}.<|im_end|>
<|im_start|>user
{question}<|im_end|>
<|im_start|>assistant
{answer}<|im_end|>
\end{verbatim}
\end{tcolorbox}

% ----------------------------------------------------------
\subsection{Answer Verification Prompt}\label{app:verify-prompt}

To verify whether the student's predicted answer matches the gold answer, we use the following grading prompt:

\begin{tcolorbox}[colback=red!3, colframe=red!40, title=Answer Verification (OREAL Correctness Judge), breakable]
\begin{verbatim}
You are a helpful assistant who evaluates the correctness 
and quality of models' outputs.
Please as a grading expert, judge whether the final answers 
given by the candidates below are consistent with the 
standard answers, that is, whether the candidates answered 
correctly.

Here are some evaluation criteria:
1. Please refer to the given standard answer. You don't need 
   to re-generate the answer to the question because the 
   standard answer has been given. You only need to judge 
   whether the candidate's answer is consistent with the 
   standard answer according to the form of the question. 
   Don't try to answer the original question. You can assume 
   that the standard answer is definitely correct.
2. Because the candidate's answer may be different from the 
   standard answer in the form of expression, before making 
   a judgment, please understand the question and the 
   standard answer first, and then judge whether the 
   candidate's answer is correct, but be careful not to try 
   to answer the original question.
3. Some answers may contain multiple items, such as 
   multiple-choice questions, multiple-select questions, 
   fill-in-the-blank questions, etc. As long as the answer 
   is the same as the standard answer, it is enough. For 
   multiple-select questions and multiple-blank 
   fill-in-the-blank questions, the candidate needs to 
   answer all the corresponding options or blanks correctly 
   to be considered correct.
4. Some answers may be expressed in different ways, such as 
   some answers may be a mathematical expression, some 
   answers may be a textual description, as long as the 
   meaning expressed is the same. And some formulas are 
   expressed in different ways, but they are equivalent 
   and correct.
5. If the prediction is given with \boxed{}, please ignore 
   the \boxed{} and only judge whether the candidate's 
   answer is consistent with the standard answer.

Please judge whether the following answers are consistent 
with the standard answer based on the above criteria. 
Grade the predicted answer of this new question as one of:
A: CORRECT
B: INCORRECT
Just return the letters "A" or "B", with no text around it.

Here is your task. Simply reply with either CORRECT, 
INCORRECT. Don't apologize or correct yourself if there 
was a mistake; we are just trying to grade the answer.

<Original Question Begin>:
{question}
<Original Question End>

<Gold Target Begin>:
{gold_answer}
<Gold Target End>

<Predicted Answer Begin>:
{answer}
<Predicted End>

Judging the correctness of candidates' answers:
\end{verbatim}
\end{tcolorbox}

% ----------------------------------------------------------
\subsection{Summary of Prompt Usage}\label{app:prompt-summary}

Table~\ref{tab:prompt-summary} summarizes the prompts and their usage across the CELPO pipeline stages.

\begin{table}[htbp]
\centering
\caption{Summary of prompts used in the CELPO framework.}
\label{tab:prompt-summary}
\small
\begin{tabular}{lll}
\toprule
\textbf{Prompt} & \textbf{Stage} & \textbf{Purpose} \\
\midrule
System Prompt & All stages & Activate step-by-step reasoning \\
Student Exam & Examination & Generate student solutions \\
Hint-Enhanced Exam & Re-examination & Prefix-force with hints \\
Teacher Correction & Hint Generation & Identify knowledge gaps \\
Training Target (Mode B) & IRDCL Training & Learn hint + answer generation \\
Training Target (Mode A) & Anchor Training & Maintain base capability \\
Answer Verification & Evaluation & Judge answer correctness \\
\bottomrule
\end{tabular}
\end{table}

% or copy the content into your appendix section.
% Requires packages: listings, xcolor, amsmath, tcolorbox (optional)
% ============================================================

\appendix
\section{Prompts}\label{app:prompts}

This appendix provides the complete prompt templates used across the different stages of our CELPO framework: inference (examination), teacher correction (hint generation), training, and answer verification.

% ----------------------------------------------------------
\subsection{System Prompt}\label{app:system-prompt}

All inference and training stages share the same system prompt to activate the model's mathematical reasoning capability:

\begin{tcolorbox}[colback=gray!5, colframe=gray!50, title=System Prompt]
\texttt{Please reason step by step and put your final answer within \textbackslash boxed\{\}.}
\end{tcolorbox}

% ----------------------------------------------------------
\subsection{Student Examination Prompt}\label{app:exam-prompt}

During the examination phase, the student model receives a math problem via the standard chat template. The input is structured as:

\begin{tcolorbox}[colback=blue!3, colframe=blue!40, title=Student Exam Input Format]
\begin{verbatim}
<|im_start|>system
Please reason step by step and put your final 
answer within \boxed{}.<|im_end|>
<|im_start|>user
{question}<|im_end|>
<|im_start|>assistant
\end{verbatim}
\end{tcolorbox}

\noindent The model then generates its reasoning and final answer in a single pass.

% ----------------------------------------------------------
\subsection{Hint-Enhanced Examination Prompt}\label{app:hint-exam-prompt}

When evaluating the student model with teacher-provided hints (prefix-forcing), the hint is prepended to the assistant's response as a prefix:

\begin{tcolorbox}[colback=blue!3, colframe=blue!40, title=Hint-Enhanced Exam Input Format]
\begin{verbatim}
<|im_start|>system
Please reason step by step and put your final 
answer within \boxed{}.<|im_end|>
<|im_start|>user
{question}<|im_end|>
<|im_start|>assistant
{hint}
\end{verbatim}
\end{tcolorbox}

\noindent The model is forced to continue generation after the hint prefix, producing the answer portion conditioned on the provided hint.

% ----------------------------------------------------------
\subsection{Teacher Correction Prompt}\label{app:teacher-correct-prompt}

The teacher model receives the following prompt to generate targeted knowledge hints based on the student's incorrect solution:

\begin{tcolorbox}[colback=green!3, colframe=green!40, title=Teacher Correction Prompt, breakable]
\begin{verbatim}
**Role:** Heuristic Logic Mentor & Knowledge Bridge

**Task:**
Your task is to analyze the **Student's Solution** in 
comparison to the **Reference Answer**. Instead of acting 
as a simple grader, you must identify the specific 
**missing logical link** or **knowledge gap** that prevents 
the student from reaching the correct conclusion.

**Goal:**
Provide the student with the necessary "Knowledge Hints" 
so that, based on their existing correct reasoning, they 
can bridge the gap and solve the problem themselves.

**Critical Constraints for Hints:**
1. **Universal Knowledge Only:** Hints must be provided as 
   **general problem-solving methods, formulas, theorems, 
   or definitions** (e.g., "Recall the formula for the area 
   of a circle: S = pi*r^2" or "Use the derivatives of 
   trigonometric functions: (sin x)' = cos x").
2. **No Specific Calculations:** Do not calculate the result 
   for the specific numbers in the problem. Provide the 
   *tool*, not the *answer*.
3. **Targeted:** The hint must directly address the specific 
   step where the student's logic broke down or stopped.
4. **Token Limit:** Each individual hint must be strictly 
   limited to **50 tokens maximum**. Keep descriptions 
   concise and focus purely on the core principle.
5. **Unidentifiable Gap:** If it is impossible to determine 
   if the student lacks specific universal knowledge (e.g., 
   the error is purely a calculation/arithmetic mistake, a 
   typo, the answer is empty/irrelevant, or the reasoning 
   is too unclear to pinpoint a missing theorem), 
   return `null`.

# Problem:
{problem}

# Student's Answer:
{student_answer}

# Reference Answer:
{ref_solution}

**Output Format:**

**Condition A:** If Constraint #5 applies (unable to identify 
a specific general knowledge gap), output strictly:
`null`

**Condition B:** If a knowledge gap is identified, respond in 
the following XML format:
<hints>
[Provide the general formula, theorem, or principle. 
 Use LaTeX for math expressions.]
...
</hints>

**Example:**
* **Student's Error:** Calculated probability as P(A)+P(B) 
  but events were not mutually exclusive.
* **Your Output:**
<hints>
For any two events A and B, the probability of their union 
is P(A ∪ B) = P(A) + P(B) - P(A ∩ B).
</hints>
\end{verbatim}
\end{tcolorbox}

% ----------------------------------------------------------
\subsection{Training Target Format (IRDCL)}\label{app:training-target}

During CELPO training, the model learns to generate both hints and answers. The training target for Mode~B (hint-augmented generation) concatenates the hint and answer:

\begin{tcolorbox}[colback=orange!3, colframe=orange!40, title=Training Target Format (Mode B)]
\begin{verbatim}
{hints}{answer}
\end{verbatim}
\end{tcolorbox}

\noindent where \texttt{\{hints\}} is the teacher-generated knowledge hint and \texttt{\{answer\}} is the correct reasoning chain ending with \verb|\boxed{...}|. The loss is computed separately on the hint tokens (with Forward KL regularization) and the answer tokens (with Reverse KL regularization).

For anchor samples (Mode~A, pure SFT), only the answer portion is used as the training target:

\begin{tcolorbox}[colback=orange!3, colframe=orange!40, title=Training Target Format (Mode A / Anchor)]
\begin{verbatim}
<|im_start|>system
Please reason step by step and put your final 
answer within \boxed{}.<|im_end|>
<|im_start|>user
{question}<|im_end|>
<|im_start|>assistant
{answer}<|im_end|>
\end{verbatim}
\end{tcolorbox}

% ----------------------------------------------------------
\subsection{Answer Verification Prompt}\label{app:verify-prompt}

To verify whether the student's predicted answer matches the gold answer, we use the following grading prompt:

\begin{tcolorbox}[colback=red!3, colframe=red!40, title=Answer Verification (OREAL Correctness Judge), breakable]
\begin{verbatim}
You are a helpful assistant who evaluates the correctness 
and quality of models' outputs.
Please as a grading expert, judge whether the final answers 
given by the candidates below are consistent with the 
standard answers, that is, whether the candidates answered 
correctly.

Here are some evaluation criteria:
1. Please refer to the given standard answer. You don't need 
   to re-generate the answer to the question because the 
   standard answer has been given. You only need to judge 
   whether the candidate's answer is consistent with the 
   standard answer according to the form of the question. 
   Don't try to answer the original question. You can assume 
   that the standard answer is definitely correct.
2. Because the candidate's answer may be different from the 
   standard answer in the form of expression, before making 
   a judgment, please understand the question and the 
   standard answer first, and then judge whether the 
   candidate's answer is correct, but be careful not to try 
   to answer the original question.
3. Some answers may contain multiple items, such as 
   multiple-choice questions, multiple-select questions, 
   fill-in-the-blank questions, etc. As long as the answer 
   is the same as the standard answer, it is enough. For 
   multiple-select questions and multiple-blank 
   fill-in-the-blank questions, the candidate needs to 
   answer all the corresponding options or blanks correctly 
   to be considered correct.
4. Some answers may be expressed in different ways, such as 
   some answers may be a mathematical expression, some 
   answers may be a textual description, as long as the 
   meaning expressed is the same. And some formulas are 
   expressed in different ways, but they are equivalent 
   and correct.
5. If the prediction is given with \boxed{}, please ignore 
   the \boxed{} and only judge whether the candidate's 
   answer is consistent with the standard answer.

Please judge whether the following answers are consistent 
with the standard answer based on the above criteria. 
Grade the predicted answer of this new question as one of:
A: CORRECT
B: INCORRECT
Just return the letters "A" or "B", with no text around it.

Here is your task. Simply reply with either CORRECT, 
INCORRECT. Don't apologize or correct yourself if there 
was a mistake; we are just trying to grade the answer.

<Original Question Begin>:
{question}
<Original Question End>

<Gold Target Begin>:
{gold_answer}
<Gold Target End>

<Predicted Answer Begin>:
{answer}
<Predicted End>

Judging the correctness of candidates' answers:
\end{verbatim}
\end{tcolorbox}

% ----------------------------------------------------------
\subsection{Summary of Prompt Usage}\label{app:prompt-summary}

Table~\ref{tab:prompt-summary} summarizes the prompts and their usage across the CELPO pipeline stages.

\begin{table}[htbp]
\centering
\caption{Summary of prompts used in the CELPO framework.}
\label{tab:prompt-summary}
\small
\begin{tabular}{lll}
\toprule
\textbf{Prompt} & \textbf{Stage} & \textbf{Purpose} \\
\midrule
System Prompt & All stages & Activate step-by-step reasoning \\
Student Exam & Examination & Generate student solutions \\
Hint-Enhanced Exam & Re-examination & Prefix-force with hints \\
Teacher Correction & Hint Generation & Identify knowledge gaps \\
Training Target (Mode B) & IRDCL Training & Learn hint + answer generation \\
Training Target (Mode A) & Anchor Training & Maintain base capability \\
Answer Verification & Evaluation & Judge answer correctness \\
\bottomrule
\end{tabular}
\end{table}
